% Architectures Lambda (traitement par lots et couche rapide fusionnés à la
% requête) et Kappa (un seul traitement de flux, relu depuis le journal).
\documentclass[border=6pt]{standalone}
\usepackage{sd}
\begin{document}
\begin{tikzpicture}[x=1mm, y=1mm, note/.append style={font=\sffamily\small},
    b/.style={soft=#1, minimum width=28mm, minimum height=10mm, font=\sffamily\small, align=center},
    lab/.style={font=\sffamily\small, text=#1, align=center}, lab/.default=Dink]
  % ---- Lambda
  \node[pillo=Dslate, anchor=west] at (-14,40) {Lambda (Marz, 2011)};
  \node[b=Dslate] (e1) at (0,20) {événements\\{\scriptsize données brutes immuables}};
  \node[b=Dblue] (lot) at (50,30) {couche lots\\{\scriptsize tout recalculer, toutes les quelques heures}};
  \node[b=Dorange] (rap) at (50,10) {couche rapide\\{\scriptsize flux, dernières heures}};
  \node[b=Dblue] (vl) at (100,30) {vues par lots};
  \node[b=Dorange] (vr) at (100,10) {vues temps réel};
  \node[b=Dteal] (q1) at (148,20) {requête\\{\scriptsize fusionne les deux}};
  \draw[arr=Dslate] (e1) -- (lot); \draw[arr=Dslate] (e1) -- (rap);
  \draw[arr=Dblue] (lot) -- (vl); \draw[arr=Dorange] (rap) -- (vr);
  \draw[arr=Dteal] (vl) -- (q1); \draw[arr=Dteal] (vr) -- (q1);
  \node[lab=Dred, anchor=north] at (50,3) {la même logique écrite deux fois, dans deux systèmes};
  % ---- Kappa
  \begin{scope}[yshift=-58mm]
    \node[pillo=Dslate, anchor=west] at (-14,40) {Kappa (Kreps, 2014)};
    \node[b=Dslate] (e2) at (0,20) {journal Kafka\\{\scriptsize rétention longue}};
    \node[b=Dblue] (j1) at (50,30) {tâche de flux v1};
    \node[b=Dteal] (t1) at (100,30) {table v1\\{\scriptsize servie}};
    \node[b=Dorange] (j2) at (50,10) {tâche de flux v2\\{\scriptsize relit depuis le début}};
    \node[b=Dorange] (t2) at (100,10) {table v2};
    \node[b=Dteal] (q2) at (148,20) {requête\\{\scriptsize lit v1, puis v2}};
    \draw[arr=Dslate] (e2) -- (j1); \draw[arr=Dslate] (e2) -- (j2);
    \draw[arr=Dblue] (j1) -- (t1); \draw[arr=Dorange] (j2) -- (t2);
    \draw[arr=Dteal] (t1) -- (q2); \draw[arr=Dteal, dash pattern=on 3pt off 2pt] (t2) -- (q2);
    \node[lab=Dslate, anchor=north, align=center] at (75,3) {une seule logique ; la v2 rattrape le journal, on bascule, on supprime la v1\\(relecture mesurée : 2 millions de clics en 10,6 à 12,2 s, un seul cœur Flink)};
  \end{scope}
\end{tikzpicture}
\end{document}
